\BourbakiExercisesSection

\begin{enumerate}
\def\labelenumi{\arabic{enumi}.}
\item
  Let G be a group and H a normal subgroup of G such that G/H is cyclic of finite order n.~Let x be an element of G whose image \(\bar{x}\) in G/H generates G/H. Let \(\phi\) be the automorphism \(h \mapsto xhx^{-1}\) of H and let \(y = x^{n}\) ; then \(y \in H\) . Show that \(\phi(y) = y\) and that \(\phi^{n}\) is the inner automorphism of H defined by y. Let \(\tau\) be the operation of Z on H defined by \((m, h) \mapsto \phi^{m}(h)\) and let \(E = H \times_{\tau} Z\) be the corresponding semi-direct product. Show that the element \((y^{-1}, n)\) of E generates a central subgroup \(C_{y}\) of E and that the quotient \(E/C_{y}\) is isomorphic to G.
\item
  Show that every central extension of Z is trivial.
\item
  Define extensions (cf.~§ 5, Exercise 10)
\end{enumerate}

\[
\mathbf {Z} / 2 \mathbf {Z} \times \mathbf {Z} / 2 \mathbf {Z} \rightarrow \mathfrak {A} _ {4} \rightarrow \mathfrak {A} _ {3}
\]

\[
\mathbf {Z} / 2 \mathbf {Z} \times \mathbf {Z} / 2 \mathbf {Z} \rightarrow \mathfrak {S} _ {4} \rightarrow \mathfrak {S} _ {3}.
\]

Show that these extensions are non-trivial, that the first admits a section and that the second does not.

¶ 4. (a) Let G be a finite group of order mn such that there exists a normal subgroup H of G which is cyclic of order m, the quotient group G/H being cyclic of order n.~Show that G is generated by two elements a, b such that \(a^{m} = e\) , \(b^{n} = a^{r}\) , \(bab^{-1} = a^{s}\) , where r and s are two integers such that \(r(s - 1)\) and \(s^{n} - 1\) are multiples of m (take a to be an element generating H, b an element of a coset generating G/H; express the elements \(b^{h}a^{k}b^{-h}\) as powers of a and apply this in particular to the cases h = n, k = 1 and h = 1, k = r).

\begin{enumerate}
\def\labelenumi{(\alph{enumi})}
\setcounter{enumi}{1}
\tightlist
\item
  *Conversely, let G(m, n, r, s) be the group defined by the presentation
\end{enumerate}

\[
\langle a, b; a ^ {m} = e, b ^ {n} = a ^ {r}, b a b ^ {- 1} = a ^ {s} \rangle ,
\]

where \(m\) and \(n\) are two integers \(\geqslant 0\) and \(r\) and \(s\) arbitrary integers (cf.~§ 7, no. 6); show that, if \(m, r(s - 1)\) and \(s^n - 1\) are not all zero, \(G(m, n, r, s)\) is a finite group of order \(qn\) , where \(q\) is the greatest common divisor of \(m, |r(s - 1)|\) and \(|s^n - 1|\) ; in this group, the subgroup H generated by \(a\) is a normal subgroup of order \(q\) and G/H is a cyclic group of order \(n\) (prove that every element of \(G(m, n, r, s)\) can be written in the form \(a^x b^y\) , where \(x\) and \(y\) are two integers such that \(0 \leqslant x \leqslant q - 1\) , \(0 \leqslant y \leqslant n - 1\) , and \(G(m, n, r, s)\) is isomorphic to the group consisting of the ordered pairs \((x, y)\) of integers subjected to the above conditions, with law of composition

\[
(x, y). (x ^ {\prime}, y ^ {\prime}) = \left\{ \begin{array}{l l} (x + x ^ {\prime} s ^ {y}, y + y ^ {\prime}) & \text { if } y + y ^ {\prime} \leqslant n - 1 \\ (x + x ^ {\prime} s ^ {y} + r, y + y ^ {\prime} - n) & \text { if } y + y ^ {\prime} \geqslant n \end{array} \right.
\]

the first coordinate of the right-hand side being a sum modulo \(q\) . Examine the cases where \(m = r(s - 1) = s^n - 1 = 0\) .

The group G(n, 2, 0, -1) is called the dihedral group of order 2n and is denoted by \(\mathbf{D}_{n}\) ; the group G(4, 2, 2, -1) is a group of order 8 called the quaternionic group and denoted by \(\mathfrak{Q}\) . Show that in \(\mathfrak{Q}\) every subgroup is normal and that the intersection of the subgroups distinct from \(\{e\}\) is a subgroup distinct from \(\{e\}\) . Prove that \(\mathbf{D}_{4}\) is not isomorphic to \(\mathfrak{Q}.\) *

\begin{enumerate}
\def\labelenumi{\arabic{enumi}.}
\setcounter{enumi}{4}
\tightlist
\item
  Let F be a group whose centre is \(\{e\}\) and A its automorphism group. F is identified with a subgroup of A by means of the homomorphism Int: \(F \rightarrow A\) ; let \(\Gamma = A/F\) .
\end{enumerate}

\begin{enumerate}
\def\labelenumi{(\alph{enumi})}
\item
  Show that the extension \(\mathbf{F} \to \mathbf{A} \to \Gamma\) is trivial only if \(\Gamma = \{e\}\) (note that the centralizer of \(\mathbf{F}\) in \(\mathbf{A}\) is \(\{e\}\) ).
\item
  Suppose that \(\Gamma = \{e\}\) . Show that every extension
\end{enumerate}

\[
\mathbf {F} \rightarrow \mathbf {E} \rightarrow \mathbf {G}
\]

of a group G by the group F is trivial.

\begin{enumerate}
\def\labelenumi{\arabic{enumi}.}
\setcounter{enumi}{5}
\tightlist
\item
  Let I be a set; write \(\mathbf{F} = \mathbf{Z}, \mathbf{E} = \mathbf{Z} \times (\mathbf{Z}/2\mathbf{Z})^{\mathrm{I}}\) and \(\mathbf{G} = \mathbf{Z}/2\mathbf{Z} \times (\mathbf{Z}/2\mathbf{Z})^{\mathrm{I}}\) .
\end{enumerate}

\begin{enumerate}
\def\labelenumi{(\alph{enumi})}
\item
  Define a non-trivial extension \(\mathbf{F} \xrightarrow{i} \mathbf{E} \xrightarrow{p} \mathbf{G}\) .
\item
  Show that, if I is infinite, E is isomorphic to \(F \times G\) .
\end{enumerate}

\begin{enumerate}
\def\labelenumi{\arabic{enumi}.}
\setcounter{enumi}{6}
\tightlist
\item
  Let G and A be two groups, A being commutative. Let \(\tau\) be a homomorphism of G into the automorphism group \(\operatorname{Aut}(\mathbf{A})\) of A; if \(g \in \mathbf{G}\) , \(a \in \mathbf{A}\) , we write \(^g a = \tau(g)(a)\) . A crossed homomorphism of G into A is any mapping \(\phi: \mathbf{G} \to \mathbf{A}\) such that
\end{enumerate}

\[
\phi (g g ^ {\prime}) = \phi (g) + ^ {g} \phi (g ^ {\prime}),
\]

the group A being written additively. The crossed homomorphisms of G into A form a group under addition denoted by Z(G, A).

\begin{enumerate}
\def\labelenumi{(\alph{enumi})}
\item
  If \(a \in \mathbf{A}\) , the mapping \(g \mapsto {}^g a - a\) is denoted by \(\theta_a\) . Show that \(a \mapsto \theta_a\) is a homomorphism \(\theta\) of \(\mathbf{A}\) into \(Z(\mathbf{G}, \mathbf{A})\) . The kernel of \(\theta\) is the subgroup of \(\mathbf{A}\) consisting of the elements invariant under \(G\) . The image of \(\theta\) is denoted by \(B(\mathbf{G}, A)\) .
\item
  Let \(\mathbf{X} = \mathbf{A} \times_{\tau} \mathbf{G}\) be the semi-direct product of \(\mathbf{G}\) by \(\mathbf{A}\) . If \(\phi\) is a mapping of \(\mathbf{G}\) into \(\mathbf{A}\) , show that \(g \mapsto (\phi(g), g)\) is a section of \(\mathbf{X}\) if and only if \(\phi\) belongs to \(\mathbf{Z}(\mathbf{G}, \mathbf{A})\) . For the sections corresponding to \(\phi_1, \phi_2 \in \mathbf{Z}(\mathbf{G}, \mathbf{A})\) to be conjugate by an element of \(\mathbf{A}\) , it is necessary and sufficient that \(\phi_1 \equiv \phi_2 \bmod. \mathbf{B}(\mathbf{G}, \mathbf{A})\) .
\item
  Let \(\mathbf{F} \xrightarrow{i} \mathbf{E} \xrightarrow{p} \mathbf{G}\) be an extension of \(\mathbf{G}\) by a group \(\mathbf{F}\) whose centre is equal to \(\mathbf{A}\) . Suppose that, if \(y \in \mathbf{E}\) and \(x = p(y)\) , then
\end{enumerate}

\[
i (x f) = y. i (f). y ^ {- 1}
\]

for all \(f \in \mathbf{A}\) . If \(\phi \in \mathbf{Z}(\mathbf{G}, \mathbf{A})\) , let \(u_{\phi}\) be the mapping of \(\mathbf{E}\) into \(\mathbf{E}\) given by the formula \(u_{\phi}(y) = i(\phi(p(y))) .y\) . Show that \(\phi \mapsto u_{\phi}\) is an isomorphism of \(\mathbf{Z}(\mathbf{G}, \mathbf{A})\) onto the automorphism group \(\operatorname{Aut}(\mathbf{E})\) of the extension \(\mathbf{E}\) . This isomorphism maps \(\mathbf{B}(\mathbf{G}, \mathbf{A})\) to the group of automorphisms \(\operatorname{Int}_{\mathbf{E}}(x)\) , where \(x\) runs through \(\mathbf{A}\) .

¶ 8. With the notation of the above exercise, C(G, A) is used to denote the group of mappings of G into A. G operates on C(G, A) by the formula

\[
\left(^ {g} \phi\right) \left(g ^ {\prime}\right) = ^ {g} \left(\phi \left(g ^ {- 1} g ^ {\prime}\right)\right).
\]

If \(\sigma\) denotes the corresponding homomorphism of \(\mathbf{G}\) into \(\operatorname{Aut}(\mathbf{C}(\mathbf{G},\mathbf{A}))\) , the semi-direct product \(\mathbf{C}(\mathbf{G},\mathbf{A}) \times_{\sigma} \mathbf{G}\) is denoted by \(\mathbf{E}_0\) .

\begin{enumerate}
\def\labelenumi{(\alph{enumi})}
\item
  Let \(\varepsilon: \mathbf{A} \to \mathbf{C}(\mathbf{G}, \mathbf{A})\) be the mapping which associates with each \(a \in \mathbf{A}\) the constant mapping equal to \(a\) . Show that \(\varepsilon\) is an injective homomorphism, compatible with the action of \(\mathbf{G}\) .
\item
  Let \(\mathbf{A} \xrightarrow{i} \mathbf{E} \xrightarrow{p} \mathbf{G}\) be an extension of \(\mathbf{G}\) by \(\mathbf{A}\) such that \(i(ga) = xi(a)x^{-1}\) if \(a \in \mathbf{A}, x \in \mathbf{E}\) and \(g = p(x)\) . Let \(\rho\) be a mapping \(\mathbf{G} \to \mathbf{E}\) such that \(p \circ \rho = \operatorname{Id}_{\mathbf{G}}\) . For all \(x \in \mathbf{E}\) , let \(\phi_x\) be the mapping of \(\mathbf{G}\) into \(\mathbf{A}\) such that
\end{enumerate}

\[
x \rho (p (x ^ {- 1}) g) = i (\phi_ {x} (g)) \rho (g) \quad \text { for   all } g \in G.
\]

Show that

\[
\phi_ {x y} = \phi_ {x} + ^ {p (x)} \phi_ {y} \quad \text { if } x, y \in \mathbf {E}.
\]

Deduce that the mapping \(\Phi: \mathbf{E} \to \mathbf{E}_0\) defined by

\[
\Phi (x) = (\phi_ {x}, p (x))
\]

is a homomorphism which makes the following diagram commutative:

\[
\begin{array}{c c c c c} \mathrm{A} & \stackrel {{i}} {{\longrightarrow}} & \mathrm{E} & \stackrel {{p}} {{\longrightarrow}} & \mathrm{G} \\ \Big \downarrow_ {\varepsilon} & & \Big \downarrow_ {\Phi} & & \Big \downarrow_ {\mathrm{Id} _ {G}} \\ \mathrm{C(G,A)} & \longrightarrow & \mathrm{E} _ {0} & \longrightarrow & \mathrm{G} \end{array} .
\]

\begin{enumerate}
\def\labelenumi{(\alph{enumi})}
\setcounter{enumi}{2}
\tightlist
\item
  A G-mean on A is a homomorphism
\end{enumerate}

\[
m \colon \mathbf {C} (\mathbf {G}, \mathbf {A}) \rightarrow \mathbf {A}
\]

satisfying the two following conditions:

\[
(c _ {1}) m \circ \varepsilon = \mathrm{Id} _ {\mathrm{A}}
\]

\[
(c _ {2}) m (^ {g} \phi) = ^ {g} m (\phi) \quad \text { if } g \in \mathrm{G}, \phi \in \mathrm{C} (\mathrm{G}, \mathrm{A}).
\]

Let \(m\) be a G-mean on A. Show that, if \(\phi\) is a crossed homomorphism and \(a = m(\phi)\) , then \(\phi = \theta_{-a}\) (cf.~Exercise 7). In particular,

\[
\mathbf {Z} (\mathbf {G}, \mathbf {A}) = \mathbf {B} (\mathbf {G}, \mathbf {A}).
\]

Show that the mapping \((\phi, g) \mapsto (m(\phi), g)\) is a homomorphism of \(\mathbf{E}_0\) into \(\mathbf{A} \times_{\tau} \mathbf{G}\) . Deduce that every extension \(\mathbf{E}\) of \(\mathbf{G}\) by \(\mathbf{A}\) satisfying the condition in \((b)\) is isomorphic to \(\mathbf{A} \times_{\tau} \mathbf{G}\) (use the composition \(\mathbf{E} \xrightarrow{\Phi} \mathbf{E}_0 \to \mathbf{A} \times_{\tau} \mathbf{G}\) ) and hence admits a section. Show that two such sections are transformed one into the other by an inner automorphism of \(\mathbf{E}\) defined by an element of \(\mathbf{A}\) (use Exercise 7 and the fact that \(\mathbf{Z}(\mathbf{G}, \mathbf{A}) = \mathbf{B}(\mathbf{G}, \mathbf{A})\) ).

\begin{enumerate}
\def\labelenumi{(\alph{enumi})}
\setcounter{enumi}{3}
\tightlist
\item
  Suppose that G is finite of order n and that the mapping \(\alpha \mapsto n\alpha\) is an automorphism of A. For all \(\phi \in \mathrm{C}(\mathrm{G}, \mathrm{A})\) , let \(m(\phi)\) denote the unique element of A such that
\end{enumerate}

\[
n. m (\phi) = \sum_ {g \in G} \phi (g).
\]

Show that m is a G-mean on A.

This applies in particular when A is finite of order prime to n.

¶ 9. Let F → E → G be an extension of finite groups. Suppose no prime number divides both the order of F and the order of G.

\begin{enumerate}
\def\labelenumi{(\alph{enumi})}
\item
  Suppose F is solvable. Show that there exists a section \(s: G \rightarrow E\) and that two such sections are conjugate by an element of F. (Argue by induction on the solvability class of F; where F is commutative, use the preceding Exercise.)
\item
  Show the existence of a section† s: G → E without assuming that F is solvable. (Argue by induction on the order of G. If p divides the order of F, choose a Sylow p-subgroup P of F and consider its normalizer N in E. The image of N in G is equal to G, cf.~Exercise 25. If N ≠ E, conclude by means of the induction hypothesis; if N = E, use (a) and the induction hypothesis applied to F/P → E/P → G.)
\end{enumerate}

¶ 10. Let G be a solvable finite group of order mn, where m and n have no common prime factor. Show that there exists a subgroup H of G of order m and that every subgroup of G whose order divides m is contained in a conjugate of H (``Hall's Theorem'').

(Argue by induction on the order of G. If \(mn \neq 1\) , choose a commutative subgroup A of G which is normal and not equal to \(\{e\}\) and whose order a is a power of a prime number. Apply this induction hypothesis to G/A and to \((m/a, n)\) or \((m, n/a)\) according to whether a divides m or n.~In the second case, use Exercise 9 to pass from G/A to G.)

\begin{enumerate}
\def\labelenumi{\arabic{enumi}.}
\setcounter{enumi}{10}
\tightlist
\item
  Show that the simple group \(A_{5}\) of order 60 contains no subgroup of order 15.
\end{enumerate}

¶ 12. Let G be a nilpotent group and H the set of elements of G of finite order.

\begin{enumerate}
\def\labelenumi{(\alph{enumi})}
\item
  Show that H is a subgroup of G.
\item
  Show that every finite subset of H generates a finite subgroup.
\item
  Show that two elements of H of relatively prime orders commute.
\end{enumerate}

\begin{enumerate}
\def\labelenumi{\arabic{enumi}.}
\setcounter{enumi}{12}
\item
  Let G be a nilpotent group. Show that G is finite (resp. countable) if and only if G/D(G) is.
\item
  Let G be a finite group. Suppose that, for all \(x, y \in G\) , the subgroup of G generated by \(\{x, y\}\) is nilpotent. Show that G is nilpotent (apply the hypothesis to x and y of orders \(p_{1}^{n_{1}}\) and \(p_{2}^{n_{2}}\) , where \(p_{1}\) and \(p_{2}\) are distinct prime numbers).
\item
  Let G be a group and X a subset of G generating G.
\end{enumerate}

\begin{enumerate}
\def\labelenumi{(\alph{enumi})}
\item
  Write \(X_{1} = X\) ; for \(n \geqslant 2\) , define \(X_{n}\) , by induction on n, as the set of commutators \((x, y)\) , \(x \in X\) , \(y \in X_{n-1}\) . Show that \(X_{n} = \{e\}\) if and only if G is nilpotent of class \(\leqslant n\) . (Argue by induction on n.~If \(X_{n} = \{e\}\) , show that the subgroup H of G generated by \(X_{n-1}\) is contained in the centre of G and apply the induction hypothesis to G/H.)
\item
  Show that \(\mathbf{C}^n (\mathbf{G})\) is a normal subgroup of \(\mathbf{G}\) generated by \(\mathbf{X}_n\) . Deduce that the image of \(\mathbf{X}_n\) in \(\mathbf{C}^n (\mathbf{G}) / \mathbf{C}^{n + 1}(\mathbf{G})\) generates the group \(\mathbf{C}^n (\mathbf{G}) / \mathbf{C}^{n + 1}(\mathbf{G})\) .
\item
  Take \(\mathbf{G} = \mathfrak{S}_m\) ( \(m \geqslant 4\) ) and \(\mathbf{X} = \{s, t\}\) , where \(s\) is a transposition and \(t\) a cycle of order \(m\) . Show that \(\mathbf{X}_2\) does not generate the group \(\mathbf{C}^2(\mathbf{G}) = \mathfrak{A}_m\) .
\end{enumerate}

\begin{enumerate}
\def\labelenumi{\arabic{enumi}.}
\setcounter{enumi}{15}
\tightlist
\item
  *Let V be a vector space over a commutative field k and
\end{enumerate}

\[
\mathbf {W} = \bigwedge^ {2} \mathbf {V}
\]

(cf.~III, § 7, no. 1). On \(\mathbf{E} = \mathbf{V} \times \mathbf{W}\) a law of composition is defined by the formula

\[
(v _ {1}, w _ {1}). (v _ {2}, w _ {2}) = (v _ {1} + v _ {2}, w _ {1} + w _ {2} + v _ {1} \wedge v _ {2}).
\]

\begin{enumerate}
\def\labelenumi{(\alph{enumi})}
\item
  Show that this law gives \(\mathbf{E}\) a group structure, that of a central extension of \(\mathbf{V}\) by \(\mathbf{W}\) .
\item
  Suppose the characteristic of \(k\) is different from 2. Show that the derived group \(D(E)\) is the set of \((0, w)\) , \(w \in W\) ; it is a group isomorphic to \(W\) . Show that, if \(\dim(V) \geqslant 4\) , there exist elements of \(D(E)\) which are not commutators.*
\end{enumerate}

\begin{enumerate}
\def\labelenumi{\arabic{enumi}.}
\setcounter{enumi}{16}
\tightlist
\item
  Let \(G\) be a finite commutative group. Show the equivalence of the following conditions:
\end{enumerate}

\begin{enumerate}
\def\labelenumi{(\alph{enumi})}
\item
  G is cyclic.
\item
  Every Sylow subgroup of G is cyclic.
\item
  For every prime number \(p\) , the number of elements \(x \in \mathbf{G}\) such that \(x^p = e\) is \(\leqslant p\) .
\end{enumerate}

\begin{enumerate}
\def\labelenumi{\arabic{enumi}.}
\setcounter{enumi}{17}
\tightlist
\item
  Let \(n\) be an integer \(\geqslant 1\) and let \(G\) be a commutative group written additively, such that \(nx = 0\) for all \(x \in G\) . Let \(H\) be a subgroup of \(G\) and \(f\) a homomorphism of \(H\) into \(\mathbf{Z} / n\mathbf{Z}\) .
\end{enumerate}

\begin{enumerate}
\def\labelenumi{(\alph{enumi})}
\item
  Let \(x \in \mathbf{G}\) and let \(q\) be the order of the image of \(x\) in \(\mathbf{G} / \mathbf{H}\) ; then \(qx \in \mathbf{H}\) . Show that there exists \(\alpha \in \mathbf{Z} / n\mathbf{Z}\) such that \(q\alpha = f(x)\) . Deduce the existence of an extension of \(f\) to the subgroup of \(\mathbf{G}\) generated by \(\mathbf{H}\) and \(x\) .
\item
  Show that \(f\) extends to a homomorphism of \(\mathbf{G}\) into \(\mathbf{Z} / n\mathbf{Z}\) (use Zorn's Lemma).
\end{enumerate}

\begin{enumerate}
\def\labelenumi{\arabic{enumi}.}
\setcounter{enumi}{18}
\tightlist
\item
  Let G be a finite commutative group.
\end{enumerate}

\begin{enumerate}
\def\labelenumi{(\alph{enumi})}
\item
  Show that there exists an element \(x \in G\) whose order n is a multiple of the orders of the elements of G. (Use the decomposition of G as a product of p-groups.) Show that the subgroup H generated by such an element is a direct factor of G (apply the above exercise to construct a homomorphism \(G \to H\) whose restriction to H is the identity).
\item
  Show that G is a direct product of cyclic groups. (Argue by induction on Card(G) using (a).) (Cf. VII, § 4, no. 6.)
\end{enumerate}

\begin{enumerate}
\def\labelenumi{\arabic{enumi}.}
\setcounter{enumi}{19}
\tightlist
\item
  Let \(G\) be a finite commutative group written additively. Let \(x = \sum_{g \in G} g\) and let \(G_2\) be the subgroup of \(G\) consisting of the elements \(g\) such that \(2g = 0\) .
\end{enumerate}

\begin{enumerate}
\def\labelenumi{(\alph{enumi})}
\item
  Show that \(x = \sum_{g \in G_2} g\) .
\item
  Show that \(x = 0\) if \(\operatorname{Card}(\mathbf{G}_2) \neq 2\) . If \(\operatorname{Card}(\mathbf{G}_2) = 2\) , show that \(x\) is the unique non-zero element of \(\mathbf{G}_2\) .
\item
  Show that \(\operatorname{Card}(\mathbf{G}_2) = 2\) if and only if the Sylow 2-subgroup of \(\mathbf{G}\) is cyclic and \(\neq 0\) .
\item
  *Let \(p\) be a prime number. Show that
\end{enumerate}

\[
(p - 1)! \equiv - 1 \pmod {p}.
\]

(Apply (b) to the multiplicative group of the field \(\mathbf{Z} / p\mathbf{Z}\) .)

\begin{enumerate}
\def\labelenumi{\arabic{enumi}.}
\setcounter{enumi}{20}
\tightlist
\item
  Let \(p\) and \(q\) be two prime numbers such that \(p > q\) .
\end{enumerate}

\begin{enumerate}
\def\labelenumi{(\alph{enumi})}
\item
  Show that every finite group G of order \(pq\) is an extension of \(\mathbf{Z}/q\mathbf{Z}\) by \(\mathbf{Z}/p\mathbf{Z}\) (note that every Sylow \(p\) -subgroup of G is normal). If further \(p \not\equiv 1 \pmod{q}\) , show that G is cyclic.
\item
  Deduce that the centre of a group cannot be of index 69.
\item
  Show that, if \(p \equiv 1 \pmod{q}\) , there exist a group of order \(pq\) whose centre is \(\{e\}\) .
\item
  Show that every non-commutative group of order 6 is isomorphic to \(\mathfrak{S}_3\) .
\end{enumerate}

\begin{enumerate}
\def\labelenumi{\arabic{enumi}.}
\setcounter{enumi}{21}
\tightlist
\item
  Let G be a finite group of order n \textgreater{} 1 and let p be the least prime number dividing n.~Let P be a Sylow p-subgroup of G and N its normalizer. Show that, if P is cyclic, P is contained in the centre of N (show that the order of N/P is prime to the order of the automorphism group of P).
\end{enumerate}

¶ 23. Let σ be an automorphism of a group G.

\begin{enumerate}
\def\labelenumi{(\alph{enumi})}
\item
  Show the equivalence of the following conditions:
\item
  \(\sigma(x) = x\) implies \(x = e\) .
\end{enumerate}

\begin{enumerate}
\def\labelenumi{(\roman{enumi})}
\setcounter{enumi}{1}
\tightlist
\item
  The mapping \(x \mapsto x^{-1}\sigma(x)\) is injective.
\end{enumerate}

When these conditions are satisfied, \(\sigma\) is said (by an abuse of language) to be without fixed point.

\begin{enumerate}
\def\labelenumi{(\alph{enumi})}
\setcounter{enumi}{1}
\item
  Suppose G is finite and \(\sigma\) without fixed point. The mapping \(x \mapsto x^{-1}\sigma(x)\) is then bijective. If H is a normal subgroup of G stable under \(\sigma\) , show that the automorphism of G/H defined by \(\sigma\) is without fixed point.
\item
  Under the hypotheses of (b), let p be a prime number. Show that there exists a Sylow p-subgroup P of G which is stable under \(\sigma\) . (If \(P_{0}\) is a Sylow p-subgroup, there exists \(y \in G\) such that \(\sigma(P_{0}) = yP_{0}y^{-1}\) ; write \(y^{-1}\) in the form \(x^{-1}\sigma(x)\) and take \(P = xP_{0}x^{-1}\) .) Show that such a subgroup is unique and contains every p-subgroup of G which is stable under \(\sigma\) .
\item
  Suppose further that \(\sigma\) is of order 2. Show that \(\sigma(g) = g^{-1}\) for all \(g \in G\) (write \(g\) in the form \(x^{-1}\sigma(x)\) with \(x \in G\) ). Deduce that \(G\) is commutative and of odd order.
\end{enumerate}

\begin{enumerate}
\def\labelenumi{\arabic{enumi}.}
\setcounter{enumi}{23}
\item
  Let G be a finite group, H a Sylow subgroup of G and N the normalizer of H. Let \(X_{1}, X_{2}\) be two subsets of the centre of H and \(s \in G\) such that \(sX_{1}s^{-1} = X_{2}\) . Show that there exists \(n \in N\) such that \(nxn^{-1} = sxs^{-1}\) for all \(x \in X_{1}\) (apply the theorem on the conjugacy of Sylow subgroups to the centralizer of \(X_{1}\) ). Deduce that two central elements of H are conjugate in G if and only if they are so in N.
\item
  Let \(\phi: G \to G'\) be a surjective homomorphism of finite groups; let \(P\) be a Sylow \(p\) -subgroup of \(G\) .
\end{enumerate}

\begin{enumerate}
\def\labelenumi{(\alph{enumi})}
\item
  Let \(\mathbf{P}_1\) be a Sylow \(p\) -subgroup of \(\mathbf{G}\) such that \(\phi(\mathbf{P}) = \phi(\mathbf{P}_1)\) . Show that there exists an element \(x\) of the kernel of \(\phi\) such that \(\mathbf{P}_1 = x\mathbf{P}x^{-1}\) .
\item
  Let N (resp. N') be the normalizer in G (resp. G') of P (resp. \(\phi(\mathrm{P})\) ). Show that \(\phi(\mathrm{N}) = \mathrm{N}'\) . In particular, if the order of \(G'\) is not divisible by p, then \(\phi(\mathrm{P}) = \{e\}\) and \(\phi(\mathrm{N}) = \mathrm{G}'\) .
\end{enumerate}

¶ 26. Let G be a finite group. G is called supersolvable if there exists a composition series \((\mathbf{G}_{i})_{0\leqslant i\leqslant n}\) of G consisting of normal subgroups of G such that the quotients \(G_{i}/G_{i+1}\) are cyclic.

\begin{enumerate}
\def\labelenumi{(\alph{enumi})}
\item
  Show that every subgroup, every quotient group and every finite product of supersolvable groups is supersolvable.
\item
  Show that nilpotent \(\Rightarrow\) supersolvable \(\Rightarrow\) solvable and that the converse implications are false.
\item
  Suppose G is supersolvable. Show that, if \(G \neq \{1\}\) , there exists a normal subgroup C of G which is cyclic of prime order. Show that the derived group \(D(G)\) of G is nilpotent. Show that every subgroup of G distinct from G which is maximal is of prime index.
\item
  A finite group G is bicyclic if there exist two elements a, b of G such that, if \(C_{a}\) and \(C_{b}\) denote the cyclic subgroups generated by a and b respectively, then \(G = C_{a}C_{b} = C_{b}C_{a}\) . Show that every bicyclic group is supersolvable. (Reduce it to proving that there exists in G a normal cyclic subgroup \(N \neq \{1\}\) . Attention may be confined to the case where \(C_{a} \cap C_{b} = \{1\}\) ; let \(m \geqslant n > 1\) be the orders of a and b respectively; by considering the elements \(b^{-1}a^{k}\) , where \(0 \leqslant k \leqslant m - 1\) , show that there exists an integer s such that \(a^{s} \neq 1\) and \(b^{-1}a^{s}b \in C_{a}\) ; the set of \(a^{s}\) with this property is the desired cyclic subgroup N.)
\end{enumerate}

\begin{enumerate}
\def\labelenumi{\arabic{enumi}.}
\setcounter{enumi}{26}
\tightlist
\item
  Let S be a finite group. S is called a minimal simple group if S is simple and non-commutative and if every subgroup of S distinct from S is solvable.
\end{enumerate}

\begin{enumerate}
\def\labelenumi{(\alph{enumi})}
\item
  Show that the alternating group \(\mathfrak{A}_n\) is minimal simple if and only if \(n = 5\) .
\item
  Let G be a finite group. Show that, if G is not solvable, there exist two subgroups H and K of G with H normal in K such that K/H is a minimal simple group.
\end{enumerate}

¶ 28. Let G be a finite group and p a prime number. An element \(s \in G\) is called p-unipotent (resp. p-regular) if its order is a power of p (resp. not divisible by p).

\begin{enumerate}
\def\labelenumi{(\alph{enumi})}
\item
  Let \(x \in G\) . Show that there exists a unique ordered pair \((u, t)\) of elements of \(G\) satisfying the following conditions: \(u\) is \(p\) -unipotent, \(t\) is \(p\) -regular, \(x = ut = tu\) . (Consider first the case where \(G\) is the cyclic group generated by \(x\) .)
\item
  Let P be a Sylow p-group of G, C its centralizer and E the set of p-regular elements of G. Show that
\end{enumerate}

\[
\operatorname{Card} (\mathbf {E}) \equiv \operatorname{Card} (\mathbf {E} \cap \mathbf {C}) \pmod {p}.
\]

Deduce that \(\operatorname{Card}(\mathbf{E}) \not\equiv 0 (\bmod p)\) . (Argue by induction on \(\operatorname{Card}(\mathbf{G})\) and reduce it to the case where \(\mathbf{C} = \mathbf{G}\) ; then use \((a)\) to show that \(\operatorname{Card}(\mathbf{E}) = (\mathbf{G}:\mathbf{P}).\) )

¶ 29. Let G be a finite group of even order and let H be a Sylow 2-subgroup of G. For all \(s \in G\) , let \(\varepsilon(s)\) be the signature of the permutation \(x \mapsto sx\) of G.

\begin{enumerate}
\def\labelenumi{(\alph{enumi})}
\item
  Let \(s \in \mathbf{H}\) . Show that \(\varepsilon(s) = -1\) if and only if \(s\) generates \(\mathbf{H}\) .
\item
  Show that \(\varepsilon\) is surjective if and only if H is cyclic.
\item
  Suppose that H is cyclic. Show that there exists one and only one normal subgroup D of G such that G is the semi-direct product of H and D. (Argue by induction on the order of H.) Show that the normalizer N of H in G is the direct product of H and \(N \cap D\) .
\item
  Show that the order of a non-commutative simple group is either odd† or divisible by 4.
\end{enumerate}

¶ 30. Let p be a prime number, r and t integers \(\geqslant0\) , S a cyclic group of order \(p^{r+t}\) and T the subgroup of S of order \(p^{t}\) .

\begin{enumerate}
\def\labelenumi{(\alph{enumi})}
\tightlist
\item
  If S operates on a finite set E, show that
\end{enumerate}

\[
\operatorname{Card} (\mathbf {E}) \equiv \operatorname{Card} (\mathbf {E} ^ {\mathrm{T}}) \pmod {p ^ {r + 1}},
\]

where \(E^{T}\) denotes the set of elements of E invariant under T.

\begin{enumerate}
\def\labelenumi{(\alph{enumi})}
\setcounter{enumi}{1}
\tightlist
\item
  Let \(m, n \geqslant 0\) . Show that
\end{enumerate}

\[
\binom {p ^ {r + t} m} {p ^ {t} n} \equiv \binom {p ^ {r} m} {n} \pmod {p ^ {r + 1}}.
\]

(Let M be a set with m elements and E the set of subsets of S × M with \(p^{t}n\) elements. Define an operation of S on S × M such that there exists a bijection of \(E^{T}\) onto the set of subsets of (S/T) × M with n elements; apply (a).)

\begin{enumerate}
\def\labelenumi{\arabic{enumi}.}
\setcounter{enumi}{30}
\tightlist
\item
  Let G be a finite group and S a Sylow p-subgroup of G. Let r, t be integers \(\geqslant0\) such that \(\operatorname{Card}(S)=p^{r+t}\) . Let E be the set of subgroups of G of order \(p^{t}\) .
\end{enumerate}

\begin{enumerate}
\def\labelenumi{(\alph{enumi})}
\item
  Suppose that S is cyclic. Show that \(\operatorname{Card}(\mathbf{E}) \equiv 1 \pmod{p^{r+1}}\) . (Let S operate on E by conjugation and use Exercise 30.) Deduce that, if the subgroup of S of order \(p^{t}\) is not normal in G, there are at least \(1 + p^{r+1}\) Sylow p-subgroups in G.
\item
  Show that \(\operatorname{Card}(\mathbf{E}) \equiv 1 \pmod{p}\) even if S is not cyclic. (Let G operate by translation on the set F of subsets of G with \(p^{t}\) elements. Show that the elements of E give distinct orbits with \((\mathrm{G}: \mathrm{S})p^{r}\) elements and that all the other orbits have a number of elements divisible by \(p^{r+1}\) . Apply Exercise 30.)
\end{enumerate}

¶ 32. Let p be a prime number and G a p-group. Let \(\mathrm{G}^{*} = \mathrm{G}^{p}\mathrm{D}(\mathrm{G})\) be the subgroup of G generated by \(\mathrm{D}(\mathrm{G})\) and the \(x^{p}, x \in G\) .

\begin{enumerate}
\def\labelenumi{(\alph{enumi})}
\item
  Show that \(\mathbf{G}^*\) is the intersection of the kernels of the homomorphisms of \(\mathbf{G}\) into \(\mathbf{Z} / p\mathbf{Z}\) .
\item
  Show that a subset S of G generates G if and only if the image of S in \(\mathbf{G} / \mathbf{G}^*\) generates \(\mathbf{G} / \mathbf{G}^*\) . Deduce that, if \((\mathbf{G}:\mathbf{G}^{*}) = p^{n}\) , the integer \(n\) is the minimum number of elements in a subset of G which generates G; in particular, G is cyclic if and only if \(n\leqslant 1\) .
\item
  Let \(u\) be an automorphism of \(G\) of order prime to \(p\) and let \(\bar{u}\) be the corresponding automorphism of \(G / G^*\) . Show that \(\bar{u} = 1\) implies \(u = 1\) .
\end{enumerate}

¶ 33. Let G be a p-group and A its automorphism group.

\begin{enumerate}
\def\labelenumi{(\alph{enumi})}
\tightlist
\item
  Let \((\mathbf{G}_i)_{0\leqslant i\leqslant n}\) be a composition series of \(\mathbf{G}\) such that \((\mathbf{G}_i:\mathbf{G}_{i + 1}) = p\) for
\end{enumerate}

\(0 \leqslant i \leqslant n - 1\) . Let \(P\) be the subgroup of \(A\) consisting of the automorphisms \(u\) such that, for all \(i\) and all \(x \in G_i\) , \(u(x)x^{-1} \in G_{i+1}\) . Show that, if an element \(u \in P\) is of order prime to \(p\) , then \(u = 1\) (argue by induction on \(n\) ). Deduce that \(P\) is a \(p\) -group.

\begin{enumerate}
\def\labelenumi{(\alph{enumi})}
\setcounter{enumi}{1}
\tightlist
\item
  Conversely, let P be a \(p\) -subgroup of A. Show that there exists a composition series \((\mathbf{G}_i)_{0 \leqslant i \leqslant n}\) of G which is stable under P and such that \((\mathbf{G}_i: \mathbf{G}_{i+1}) = p\) for \(0 \leqslant i \leqslant n-1\) ; show that the \(\mathbf{G}_i\) may be chosen to be normal in G.
\end{enumerate}

\begin{enumerate}
\def\labelenumi{\arabic{enumi}.}
\setcounter{enumi}{33}
\item
  Let p be a prime number. Show that every group of order \(p^{2}\) is commutative.
\item
  Let G be a group such that the derived group \(D = D(G)\) is contained in the centre C of G.
\end{enumerate}

\begin{enumerate}
\def\labelenumi{(\alph{enumi})}
\tightlist
\item
  Show that the mapping \((x,y)\mapsto (x,y)\) induces a mapping
\end{enumerate}

\[
\phi \colon (\mathbf {G} / \mathbf {C}) \times (\mathbf {G} / \mathbf {C}) \rightarrow \mathbf {D}
\]

such that, writing these groups additively:

\[
\begin{array}{l l} \phi (\alpha + \beta , \gamma) = \phi (\alpha , \gamma) + \phi (\beta , \gamma) & \quad \phi (\alpha , \beta) = - \phi (\beta , \alpha) \\ \phi (\alpha , \beta + \gamma) = \phi (\alpha , \beta) + \phi (\alpha , \gamma) & \quad \phi (\alpha , \alpha) = 0 \end{array}
\]

for all \(\alpha, \beta, \gamma \in \mathbf{G} / \mathbf{C}\) .

\begin{enumerate}
\def\labelenumi{(\alph{enumi})}
\setcounter{enumi}{1}
\tightlist
\item
  Show that, for every integer n,
\end{enumerate}

\[
(y x) ^ {n} = y ^ {n} x ^ {n} (x, y) ^ {n (n - 1) / 2}
\]

for \(x, y \in \mathbf{G}\) . Deduce that if \(n\) is odd and \(d^n = 1\) for all \(d \in \mathbf{D}\) , the mapping \(x \mapsto x^n\) induces a homomorphism \(\theta: \mathbf{G} / \mathbf{D} \to \mathbf{C}\) .

¶ 36. Let p be a prime number and G a non-commutative group of order \(p^{3}\) .

\begin{enumerate}
\def\labelenumi{(\alph{enumi})}
\item
  Show that, with the notation of the preceding exercise, \(\mathbf{C} = \mathbf{D}\) and that \(\mathbf{G} / \mathbf{D}\) is isomorphic to the product of two groups of order \(p\) .
\item
  Suppose that p is odd and that the homomorphism \(\theta\) of the preceding exercise is non-zero. Show that G is the semi-direct product of a group of order p by a cyclic group of order \(p^{2}\) . Show that there exist elements x, y generating G such that
\end{enumerate}

\[
x ^ {p ^ {2}} = 1, \quad y ^ {p} = 1, \quad (x, y) = x ^ {p},
\]

that G is characterized up to isomorphism by this property (and by the fact that it is of order \(p^{3}\) ) and that such a group G exists.

\begin{enumerate}
\def\labelenumi{(\alph{enumi})}
\setcounter{enumi}{2}
\item
  For \(p = 2\) , consider the same question as in (b) with the hypothesis \(\theta \neq 0\) replaced by the hypothesis that \(G\) contains a non-central element of order 2. Show that such a group is isomorphic to the dihedral group \(D_4\) (cf.~Exercise 4) and also to a Sylow 2-subgroup of \(S_4\) .
\item
  Suppose that \(p\) is odd and that the homomorphism \(\theta\) of the preceding exercise is zero. Show that there are elements \(x, y, z\) generating \(G\) such that
\end{enumerate}

\[
(x, y) = z; \quad x ^ {p} = y ^ {p} = z ^ {p} = (x, z) = (y, z) = 1,
\]

that G is characterized by this property and that such a group exists.

*Show that G is isomorphic to the multiplicative group of matrices of the form

\[
\left( \begin{array}{c c c} 1 & a & b \\ 0 & 1 & c \\ 0 & 0 & 1 \end{array} \right)
\]

with coefficients \(a, b, c\) in the field with \(p\) elements.*

\begin{enumerate}
\def\labelenumi{(\alph{enumi})}
\setcounter{enumi}{4}
\tightlist
\item
  Suppose that \(p = 2\) and that no non-central element of \(G\) is of order 2 (hence that they are all of order 4). Show that \(G\) is generated by elements \(x, y\) such that
\end{enumerate}

\[
x ^ {2} = y ^ {2} = (x, y); \quad (x, y) ^ {2} = 1,
\]

that G is characterized by this property and is isomorphic to the quaternionic group (cf.~Exercise 4).

\begin{enumerate}
\def\labelenumi{(\alph{enumi})}
\setcounter{enumi}{5}
\tightlist
\item
  *Is the group of matrices
\end{enumerate}

\[
\left( \begin{array}{c c c} 1 & a & b \\ 0 & 1 & c \\ 0 & 0 & 1 \end{array} \right),
\]

where \(a, b, c\) run through the field with 2 elements, of type \((c)\) or of type \((e)?\) *

\begin{enumerate}
\def\labelenumi{\arabic{enumi}.}
\setcounter{enumi}{36}
\tightlist
\item
  \begin{enumerate}
  \def\labelenumii{(\alph{enumii})}
  \tightlist
  \item
    Let G be a finite group, H a normal subgroup of G, p a prime number not dividing the order of H and P a Sylow p-subgroup of G. Show that HP is the semi-direct product of P by H.
  \end{enumerate}
\end{enumerate}

\begin{enumerate}
\def\labelenumi{(\alph{enumi})}
\setcounter{enumi}{1}
\tightlist
\item
  We say that a finite group G has property (ST) if there exists a numbering \(p_{1}, \ldots, p_{s}\) of distinct prime numbers dividing the order
\end{enumerate}

\[
n = p _ {1} ^ {a _ {1}} \dots p _ {s} ^ {a _ {s}}
\]

of G and Sylow \(p_i\) -subgroups \(P_1, \ldots, P_s\) of G such that, for \(1 \leqslant i \leqslant s\) , the set \(G_i = P_1P_2 \ldots P_i\) is a normal subgroup of G. Show that, if this is so, \(G_i\) does not depend on the choice of the \(P_i\) ; it is the unique subgroup of G of order \(p_1^{a_1} \ldots p_i^{a_i}\) ; further, \(G_i\) is the semi-direct product of \(P_i\) by \(G_{i-1}\) .

\begin{enumerate}
\def\labelenumi{(\alph{enumi})}
\setcounter{enumi}{2}
\item
  Show that a finite group G has property (ST) if and only if there exists a Sylow subgroup P of G which is normal in G and such that G/P has property (ST).
\item
  Show that every subgroup, every quotient group and every central extension of a group with property (ST) has property (ST).
\end{enumerate}

¶ 38. Let G be a finite group of order \(n = p^{a}m\) , where \(p\) is a prime number not dividing \(m\) , and let X be the set of Sylow \(p\) -subgroups of G. Let P \ensuremath{\in} X, let N be the normalizer of P and \(r = \text{Card}(X)\) .

\begin{enumerate}
\def\labelenumi{(\alph{enumi})}
\item
  Show that \(r = (\mathbf{G}:\mathbf{N})\) . Deduce that \(r\) belongs to the set \(\mathbf{R}\) of positive divisors of \(m\) which are congruent to 1 (mod \(p\) ). In particular, if \(\mathbf{R}\) is \(\{1\}\) , \(\mathbf{P}\) is normal in \(\mathbf{G}\) .
\item
  If \(a = 1\) , show that the number of elements of \(G\) of order \(p\) is \(r(p - 1)\) . If further \(N = P\) , show that the elements of \(G\) of order different from \(p\) are equal in number to \(m\) ; if further \(m\) is a power of a prime number \(q\) , deduce that a Sylow \(q\) -subgroup \(Q\) of \(G\) contains all the elements of \(G\) of order different from \(p\) and hence is normal in \(G\) .
\item
  Let H be a group of order 210. Show that H contains a normal subgroup G of order 105 (cf.~Exercise 29). Let p = 3 (resp. 5, 7). Show that, if G has no normal Sylow p-subgroup, the group G contains at least 14 (resp. 84, 90) elements of order p.~Conclude that H and G have property (ST) (cf.~the preceding exercise).
\item
  Show that every group of order \(n \leqslant 100\) has property (ST), except possibly if \(n = 24, 36, 48, 60, 72, 96\) (same method as in (c)).
\item
  Show that the group \(\mathfrak{S}_4\) , of order 24, does not have property (ST). Construct analogous examples for \(n = 48, 60, 72, 96\) . (For \(n = 36\) , see the following exercise.)
\end{enumerate}

¶ 39. Let G be a finite group, \(n = \text{Card}(G)\) , \(p\) a prime number, X the set of Sylow \(p\) -subgroups of G and \(r = \text{Card}(X)\) . Then G operates transitively on X by conjugation; let \(\phi: G \to \mathfrak{S}_{X} \approx \mathfrak{S}_r\) be the corresponding homomorphism ( \(A \approx B\) denotes the relation ``A is isomorphic to B'') and K its kernel.

\begin{enumerate}
\def\labelenumi{(\alph{enumi})}
\item
  Show that K is the intersection of the normalizers of the \(P \in X\) and that the Sylow p-subgroup of K is the intersection of the \(P \in X\) . Deduce that if r \textgreater{} 1 the order k of K divides n/rp.
\item
  Show that if G has no subgroup of index 2, then \(\phi(G) \subset \mathfrak{A}_{X}\) and that if \(n = kr!\) (resp. \(\frac{1}{2}kr!\) ), then \(\phi(G) = \mathfrak{S}_{X}\) (resp. \(\mathfrak{A}_{X}\) ).
\item
  Using (a), (b) and the preceding exercise, show the following facts:
\end{enumerate}

If \(n = 12\) and a Sylow 3-subgroup of \(G\) is not normal, then \(G \approx \mathfrak{A}_4\) .

If \(n = 24\) and \(\mathbf{G}\) does not have property (ST), then \(\mathrm{G} \approx \mathfrak{S}_4\) .

If n = 36, then G has property (ST). (Show that if a Sylow 3-subgroup of G is not normal, then G contains a subgroup K of order 3 such that \(G/K \approx A_{4}\) and that K is central.)

If \(n = 48\) and \(G\) does not have property (ST), then \(G\) contains a normal subgroup \(K\) of order 2 such that \(G / K \approx S_4\) .

If \(n = 60\) and \(G\) does not have property (ST), then \(G \approx \mathfrak{A}_5\) . (Show that such a group \(G\) has no normal subgroup of order 5, nor of order 2, nor of index 2. Deduce that \(G\) is isomorphic to a subgroup of \(\mathfrak{S}_6\) and, arguing as in Exercise 25 of § 5, that \(G \approx \mathfrak{A}_5\) .)

If \(n = 72\) and a Sylow 3-subgroup of \(G\) is not normal, then \(G\) contains a normal subgroup \(K\) such that \(G / K\) is isomorphic either to \(\mathfrak{S}_4\) or to \(\mathfrak{A}_4\) .

If n = 96 and a Sylow 2-subgroup of G is not normal, then G contains a normal subgroup K such that G/K ≈ \(S_{3}\) .

\begin{enumerate}
\def\labelenumi{(\alph{enumi})}
\setcounter{enumi}{3}
\tightlist
\item
  Show that a group G of order \(\leqslant 100\) is solvable unless \(\mathbf{G} \approx \mathfrak{A}_5\) .
\end{enumerate}

\begin{enumerate}
\def\labelenumi{\arabic{enumi}.}
\setcounter{enumi}{39}
\tightlist
\item
  Let \(G_{1}, G_{2}\) be two groups and G a subgroup of \(G_{1} \times G_{2}\) such that \(pr_{1} G = G_{1}, pr_{2} G = G_{2}\) ; identifying \(G_{1}\) and \(G_{2}\) with \(G_{1} \times \{e_{2}\}\) and \(\{e_{1}\} \times G_{2}\) respectively, \(H_{1} = G \cap G_{1}\) and \(H_{2} = G \cap G_{2}\) are normal subgroups of G such that \(G/H_{1}\) is isomorphic to \(G_{2}\) and \(G/H_{2}\) isomorphic to \(G_{1}\) .
\end{enumerate}

\begin{enumerate}
\def\labelenumi{(\alph{enumi})}
\item
  Show that for G to be normal in \(G_{1} \times G_{2}\) , it is necessary and sufficient that G contain the commutator group \(\mathrm{D}(\mathrm{G}_{1} \times \mathrm{G}_{2}) = \mathrm{D}(\mathrm{G}_{1}) \times \mathrm{D}(\mathrm{G}_{2})\) (prove that G contains \(\mathrm{D}(\mathrm{G}_{1})\) by writing \(xzx^{-1}z^{-1} \in G\) for \(x \in G_{1}\) and \(z \in G\) ). Deduce that \(G_{1}/H_{1}\) is commutative.
\item
  If G is normal in \(G_{1} \times G_{2}\) , show that a homomorphism is defined of \(G_{1} \times G_{2}\) onto \(G_{1}/H_{1}\) by associating with each ordered pair \((s_{1}, s_{2})\) the coset of \(s_{1}t_{1}^{-1}\) modulo \(H_{1}\) , where \(t_{1} \in G_{1}\) is such that \((t_{1}, s_{2}) \in G\) . Deduce that \((\mathrm{G}_{1} \times \mathrm{G}_{2})/\mathrm{G}\) is isomorphic to \(\mathrm{G}/(\mathrm{H}_{1} \times \mathrm{H}_{2})\) .
\end{enumerate}

\begin{enumerate}
\def\labelenumi{\arabic{enumi}.}
\setcounter{enumi}{40}
\tightlist
\item
  Let \(x, y\) be two elements of a group \(G\) .
\end{enumerate}

\begin{enumerate}
\def\labelenumi{(\alph{enumi})}
\item
  For there to exist \(a, b\) in \(G\) such that \(bay = xab\) , it is necessary and sufficient that \(xy^{-1}\) be a commutator.
\item
  For there to exist \(2n + 1\) elements \(a_1, a_2, \ldots, a_{2n+1}\) in G such that
\end{enumerate}

\[
x = a _ {1} a _ {2} \dots a _ {2 n + 1} \quad \text { and } \quad y = a _ {2 n + 1} a _ {2 n} \dots a _ {1},
\]

it is necessary and sufficient that \(xy^{-1}\) be a product of n commutators (argue by induction on n).
% semantic-fragments: legacy-input-seam
\relax{}
